\section{Modelos y protocolo de evaluación}
\label{sec:protocolo}

\subsection{Modelos considerados}
\guia{Incluir el modelo objetivo, del orden de 1 a 2 mil millones de
parámetros según el enunciado, y los modelos usados para comparación.
El curso pide modelos de frontera, privilegiando modelos abiertos.
Añadir o quitar filas según los modelos efectivamente evaluados.}

\begin{table}[htbp]
  \centering\small
  \caption{Modelos utilizados en la línea base.}
  \label{tab:modelos}
  \begin{tabularx}{\linewidth}{@{}lYYY@{}}
    \toprule
    Rol & Modelo y variante & Parámetros & Versión o revisión \\
    \midrule
    Objetivo & \dato & \dato & \dato \\
    Comparación 1 & \dato & \dato & \dato \\
    Comparación 2 & \dato & \dato & \dato \\
    \bottomrule
  \end{tabularx}
\end{table}

\subsection{Configuración y reproducibilidad}
\begin{description}[leftmargin=0pt]
  \item[Prompt y herramientas:] \pendiente{plantilla de prompt,
  formato de salida, herramientas permitidas y límite de turnos.}
  \item[Inferencia:] \pendiente{temperatura, top-p, máximo de tokens,
  semillas y cualquier configuración que varíe entre modelos.}
  \item[Entorno de ejecución:] \pendiente{hardware, precisión o
  cuantización, software y versiones relevantes.}
  \item[Verificación:] \pendiente{versión del verificador, criterio
  binario de éxito y tratamiento de errores, salidas inválidas y
  tiempos de espera.}
\end{description}

\subsection{Métricas y muestreo}
\pendiente{definir cómo se calculan \passuno\ y \passk, el valor o
valores de $k$, el número de problemas por nivel y el número de
respuestas muestreadas por problema.}

\guia{Explicitar el estimador utilizado y citar su fuente si
corresponde. Distinguir el número total de muestras $n$ del valor
$k$; debe haber suficientes muestras para estimar la métrica elegida.
Reportar siempre la misma escala, proporción o porcentaje, e indicar
cómo se agregan los resultados. El enunciado no fija un valor de $k$.
Una recompensa parcial no equivale automáticamente a un acierto.}
